\hypertarget{python-3.1-to-3.2-standard-library-consolidation-and-threading-pools}{%
\chapter{Python 3.1 to 3.2: Standard Library Consolidation and Threading
Pools}\label{python-3.1-to-3.2-standard-library-consolidation-and-threading-pools}}

\hypertarget{antoine-pitrous-new-gil-python-3.2}{%
\subsection{8.1 Antoine Pitrou's New GIL (Python
3.2)}\label{antoine-pitrous-new-gil-python-3.2}}

To resolve multi-core performance bottlenecks, Python 3.2 replaced the
legacy ticker-based GIL with an interval-based GIL designed by Antoine
Pitrou.

\hypertarget{the-convoy-effect-gil-battle-under-the-ticker-based-gil}{%
\subsubsection{1. The Convoy Effect \& GIL Battle under the Ticker-Based
GIL}\label{the-convoy-effect-gil-battle-under-the-ticker-based-gil}}

In Python 2.x and pre-3.2, thread switching was based on a bytecode
execution ticker (\passthrough{\lstinline!sys.checkinterval!}, default
100 instructions). Once a thread \(T_1\) executed 100 bytecodes, it
released the GIL, signaled waiting threads, and immediately re-attempted
to acquire it: 1. \(T_1\) releases GIL \(\rightarrow\) \(T_1\)
immediately calls \passthrough{\lstinline!acquire!} again. 2. Because
\(T_1\) is already running on a CPU core with warm caches, it is highly
likely to re-acquire the GIL before a sleeping thread \(T_2\) on another
core can wake up. 3. This resulted in the \textbf{convoy effect} (or GIL
battle), causing thread starvation and wasting CPU cycles on rapid,
unsuccessful mutex context-switching.

\hypertarget{the-interval-based-gil-architecture}{%
\subsubsection{2. The Interval-Based GIL
Architecture}\label{the-interval-based-gil-architecture}}

The new GIL manages thread switching using a time interval (default 5000
microseconds / 5ms). The state of the GIL is maintained in a global
runtime structure (\passthrough{\lstinline!ceval\_gil.h!}):

\begin{lstlisting}[language=C]
typedef struct {
    Py_MUTEX_T mutex;              /* Mutex guarding the GIL status fields */
    Py_COND_T cond;                /* Condition variable for waiting threads */
    int locked;                    /* Flag indicating if the GIL is locked (1 or 0) */
    unsigned long switch_interval; /* Maximum thread execution duration before switch (in microseconds) */
    
    PyThreadState* volatile tstate; /* Pointer to the thread state currently holding the GIL */
    int volatile gil_drop_request;  /* Shared atomic flag indicating a thread must drop the GIL */
} _gil_runtime_state;
\end{lstlisting}

\hypertarget{step-by-step-thread-state-transitions-mutex-flow}{%
\subsubsection{3. Step-by-Step Thread State Transitions \& Mutex
Flow}\label{step-by-step-thread-state-transitions-mutex-flow}}

When thread \(T_1\) holds the GIL and thread \(T_2\) requests execution:

\begin{lstlisting}
Thread 1 (Active)                               Thread 2 (Waiting)
-----------------                               ------------------
Holds GIL (locked=1)
Executes bytecodes...                           Requests GIL via take_gil()
                                                Acquires mutex, waits on cond with 5ms timeout
                                                Timeout expires! (T1 hasn't released GIL)
                                                Sets gil_drop_request = 1
Checks gil_drop_request in eval loop
Detects flag set to 1
Releases GIL via drop_gil()
Resets locked=0, tstate=NULL
Signals cond, releases mutex
                                                Wakes up from cond, acquires mutex
                                                Sets locked=1, tstate=T2
                                                Resets gil_drop_request=0
                                                Starts executing bytecodes...
\end{lstlisting}

\begin{enumerate}
\def\labelenumi{\arabic{enumi}.}
\item
  \textbf{Requesting the GIL}: Thread \(T_2\) enters
  \passthrough{\lstinline!take\_gil()!}, acquires the
  \passthrough{\lstinline!mutex!}, and waits on the condition variable
  \passthrough{\lstinline!cond!} with a timeout of
  \passthrough{\lstinline!switch\_interval!} (5ms).
\item
  \textbf{Timeout \& Requesting Release}: If the timeout expires and
  \(T_1\) has not released the GIL (e.g., due to executing an I/O
  operation or blocking system call), \(T_2\) sets the global atomic
  flag \passthrough{\lstinline!gil\_drop\_request = 1!}.
\item
  \textbf{Interpreter Eval Loop Check}: In CPython's evaluation loop
  (\passthrough{\lstinline!\_PyEval\_EvalFrameDefault!} in
  \passthrough{\lstinline!Python/ceval.c!}), the interpreter checks
  \passthrough{\lstinline!gil\_drop\_request!} between bytecode
  execution steps:

\begin{lstlisting}[language=C]
/* Check GIL drop request flag */
if (_Py_atomic_load_relaxed(&gil_runtime_state.gil_drop_request)) {
    /* Give up the GIL and wait for reschedule */
    PyThreadState *tstate = _PyThreadState_GET();
    PyEval_SaveThread();     /* Drops GIL, signals cond */
    PyEval_RestoreThread();  /* Re-acquires GIL, blocks if held */
}
\end{lstlisting}
\item
  \textbf{Acquiring the GIL}: Once \(T_1\) calls
  \passthrough{\lstinline!PyEval\_SaveThread()!}, it sets
  \passthrough{\lstinline!locked = 0!}, signals
  \passthrough{\lstinline!cond!} to wake up \(T_2\), and releases the
  \passthrough{\lstinline!mutex!}. \(T_2\) wakes up, sets
  \passthrough{\lstinline!locked = 1!}, clears
  \passthrough{\lstinline!gil\_drop\_request = 0!}, and starts its
  execution phase.
\end{enumerate}

\begin{center}\rule{0.5\linewidth}{0.5pt}\end{center}

\hypertarget{thread-pools-and-process-pools-concurrent.futures-pep-3148}{%
\subsection{\texorpdfstring{8.2 Thread Pools and Process Pools
(\texttt{concurrent.futures} / PEP
3148)}{8.2 Thread Pools and Process Pools (concurrent.futures / PEP 3148)}}\label{thread-pools-and-process-pools-concurrent.futures-pep-3148}}

PEP 3148 unified concurrent task execution under a shared API using
\passthrough{\lstinline!ThreadPoolExecutor!} and
\passthrough{\lstinline!ProcessPoolExecutor!}.

\hypertarget{threadpoolexecutor-mechanics}{%
\subsubsection{\texorpdfstring{1. \texttt{ThreadPoolExecutor}
Mechanics}{1. ThreadPoolExecutor Mechanics}}\label{threadpoolexecutor-mechanics}}

\passthrough{\lstinline!ThreadPoolExecutor!} manages a pool of worker
threads using a task queue
(\passthrough{\lstinline!queue.SimpleQueue!}): * \textbf{Task
Submission}: When \passthrough{\lstinline!submit(fn, *args)!} is called,
CPython wraps the callable in a \passthrough{\lstinline!Future!} object
and a \passthrough{\lstinline!\_WorkItem!} struct, pushing it to the
executor's task queue. * \textbf{Worker Execution}: Worker threads
execute a loop inside \passthrough{\lstinline!\_worker()!}:
\passthrough{\lstinline!python   def \_worker(executor\_reference, work\_queue):       while True:           work\_item = work\_queue.get(block=True)           if work\_item is not None:               work\_item.run()!}
* \textbf{GIL Constraints}: Because all worker threads reside in the
same memory space, they share the GIL. They can execute I/O-bound
operations (e.g., waiting on socket descriptors or file operations)
concurrently by releasing the GIL at the C-API level during blocking
calls, but cannot execute CPU-bound tasks in parallel.

\hypertarget{processpoolexecutor-ipc-serialization-bottlenecks}{%
\subsubsection{\texorpdfstring{2. \texttt{ProcessPoolExecutor} \& IPC
Serialization
Bottlenecks}{2. ProcessPoolExecutor \& IPC Serialization Bottlenecks}}\label{processpoolexecutor-ipc-serialization-bottlenecks}}

To bypass the GIL for CPU-bound operations,
\passthrough{\lstinline!ProcessPoolExecutor!} spawns independent worker
processes. Because processes do not share memory, CPython must serialize
task arguments and deserialize results using the
\passthrough{\lstinline!pickle!} protocol.

The math of Process Pool Execution Overhead:
\[\text{Total Execution Time} = t_{\text{serialization}} + t_{\text{IPC transfer}} + t_{\text{computation}} + t_{\text{IPC return}} + t_{\text{deserialization}}\]
\[\text{IPC Overhead} \propto \text{size of arguments} + \text{size of results}\]
1. \textbf{Serialization}: The parent process pickles the function and
arguments into a byte stream. 2. \textbf{IPC Transfer}: The byte stream
is written to an OS pipe or socket descriptor, crossing the user-kernel
space boundary twice:
\[\text{Parent User Space} \xrightarrow{\text{write()}} \text{Kernel Buffer} \xrightarrow{\text{read()}} \text{Child User Space}\]
3. \textbf{Computation}: The child process runs the task within its own
interpreter instance and GIL. 4. \textbf{Return}: The child pickles and
writes the result back through a return pipe. For small, fast
computations, the serialization and pipe I/O overhead can exceed the
execution time of the computation itself.

\hypertarget{future-state-tracking}{%
\subsubsection{3. Future State Tracking}\label{future-state-tracking}}

A \passthrough{\lstinline!Future!} object tracks task completion state.
The state transitions are guarded by a thread-local reentrant lock
(\passthrough{\lstinline!threading.RLock!}):

\begin{lstlisting}[language=Python]
# Future state definitions inside concurrent.futures.futures
PENDING = 'PENDING'
RUNNING = 'RUNNING'
CANCELLED = 'CANCELLED'
FINISHED = 'FINISHED'
\end{lstlisting}

\begin{itemize}
\tightlist
\item
  Calling \passthrough{\lstinline!cancel()!} transitions the state from
  \passthrough{\lstinline!PENDING!} to
  \passthrough{\lstinline!CANCELLED!}, waking up any threads waiting on
  \passthrough{\lstinline!result()!} via a condition variable.
\item
  When the task completes, \passthrough{\lstinline!set\_result(val)!}
  transitions the state to \passthrough{\lstinline!FINISHED!}, stores
  the return value in \passthrough{\lstinline!self.\_result!}, and
  invokes all callback functions registered via
  \passthrough{\lstinline!add\_done\_callback(fn)!}.
\end{itemize}

\begin{center}\rule{0.5\linewidth}{0.5pt}\end{center}

\hypertarget{standard-library-consolidation}{%
\subsection{8.3 Standard Library
Consolidation}\label{standard-library-consolidation}}

\hypertarget{ordereddict-pep-372}{%
\subsubsection{\texorpdfstring{1. \texttt{OrderedDict} (PEP
372)}{1. OrderedDict (PEP 372)}}\label{ordereddict-pep-372}}

Introduced in Python 3.1, \passthrough{\lstinline!OrderedDict!}
preserves key insertion order. Since standard dictionaries were
unordered at this time, \passthrough{\lstinline!OrderedDict!} maintained
order by wrapping a standard dictionary with a private doubly-linked
list. The keys are stored as nodes in the linked list:

\begin{lstlisting}[language=Python]
# Doubly-linked list node format: [PREV, NEXT, KEY]
root = []
root[:] = [root, root, None] # Pointer loop representing empty list
\end{lstlisting}

\begin{itemize}
\item
  \textbf{Insertion}: When a new key is added,
  \passthrough{\lstinline!OrderedDict!} appends it to the end of the
  doubly-linked list by adjusting pointers:

\begin{lstlisting}[language=Python]
last = root[0]
last[1] = root[0] = [last, root, key]
self.__map[key] = root[0]
\end{lstlisting}
\item
  \textbf{Deletion}: Deleting a key updates the pointers of the
  neighboring nodes in \(O(1)\) time:

\begin{lstlisting}[language=Python]
link_prev, link_next, key = self.__map.pop(key)
link_prev[1] = link_next
link_next[0] = link_prev
\end{lstlisting}
\end{itemize}

\hypertarget{counter}{%
\subsubsection{\texorpdfstring{2.
\texttt{Counter}}{2. Counter}}\label{counter}}

A dictionary subclass designed for tallying hashable elements. It
optimizes standard dictionary lookups by wrapping loops in C:

\begin{lstlisting}[language=Python]
# Counter update logic
for elem in iterable:
    self[elem] = self.get(elem, 0) + 1
\end{lstlisting}

At the C level, it accesses dictionary lookups via
\passthrough{\lstinline!PyDict\_GetItemWithError!} to bypass
Python-level attribute lookup overhead, speeding up frequency counting.
\#\#\#\# 3. \passthrough{\lstinline!argparse!} Introduced in Python 3.2
to replace \passthrough{\lstinline!optparse!},
\passthrough{\lstinline!argparse!} uses a parser state machine to
process command-line arguments. It builds a hierarchical action registry
(\passthrough{\lstinline!\_actions!} list) containing
\passthrough{\lstinline!Action!} objects (e.g.,
\passthrough{\lstinline!\_StoreAction!},
\passthrough{\lstinline!\_StoreConstAction!}). During
\passthrough{\lstinline!parse\_args()!}, it iterates over argument
strings, matches positional arguments and optional flags (using a
regular expression mapping state machine), and applies type coercion
hooks (e.g., \passthrough{\lstinline!int!},
\passthrough{\lstinline!float!}) before writing variables to a
\passthrough{\lstinline!Namespace!} object.
